\section{Narrow phase}
\label{sec:np}

Each candidate pair that survives \cref{sec:bp} is decided here, by branch and
bound over boxes in the parameter domain. Two properties from \cref{sec:prelim}
make the search possible at all. \Cref{eq:hull} says that a box can be judged
from its eight corners, and \cref{eq:err} says by how much that judgement can be
wrong. \Cref{alg:np} is the skeleton every kernel in this paper instantiates.
Three things vary between them: the acceptance test, the split rule, and how the
work list is organised. The third is the subject of \cref{sec:parallel}.

\begin{algorithm}[htbp]
\caption{Branch and bound for one query. $\toi$ is the running minimum, shared
across queries when the caller asks only for the earliest impact.}
\label{alg:np}
\begin{algorithmic}[1]
\Function{NarrowPhase}{query, $\toi$}
  \State $\epsilon \gets$ error bound of \cref{eq:err};\quad
         $\tau_{0..2} \gets$ tolerances of \cref{eq:tol}
  \State push $\dom_{0} = [0,\min(\toi,1)]\times[0,1]^{2}$
  \While{a box $\dom$ remains}
    \If{$\dom.t_{\ell} \ge \toi$} \textbf{continue}
      \Comment{a root here cannot improve the answer}
    \EndIf
    \State $[\,\underline{F}_{d},\overline{F}_{d}\,] \gets$ hull of the eight
           corner values, per component \Comment{\cref{eq:hull}}
    \If{$\underline{F}_{d} > \epsilon_{d}$ or $\overline{F}_{d} < -\epsilon_{d}$
        for any $d$}
      \State \textbf{continue} \Comment{\textbf{reject}: the only unsound step,
             so it is padded by $\epsilon$}
    \EndIf
    \If{\Call{Accept}{$\dom$} \textbf{or} $\dom.\text{depth} \ge D$
        \textbf{or} $\dom$ cannot be split}
      \State $\toi \gets \min(\toi,\ \dom.t_{\ell})$;\quad \textbf{continue}
        \Comment{exhaustion \emph{accepts}, never drops}
    \EndIf
    \State push the children of \Call{Split}{$\dom$}
  \EndWhile
  \State \Return $\toi$
\EndFunction
\end{algorithmic}
\end{algorithm}

On acceptance \cref{alg:np} reports the \emph{lower} bound in $t$ of the
accepted box, never its midpoint and never its upper bound. Every exit that is
not a rejection accepts. Those two facts are the whole of the safety argument,
and \cref{sec:guarantee} makes them precise.

\subsection{Two acceptance tests}
\label{sec:np:modes}

Two instantiations of \op{Accept} and \op{Split} appear in what follows. Both
are conservative. They differ in how tight an answer they return, and in how
long they take to return it.

The first, \textsc{Tight}, is TightInclusion's test transcribed:
\begin{equation} \label{eq:accept-tight} \textsc{Accept}(\dom) \;=\;
\underbrace{\bigl(\forall d:\ [\underline{F}_{d},\overline{F}_{d}] \subseteq
[-\epsilon_{d},\epsilon_{d}]\bigr)}_{\text{function box inside the error box}}
\ \vee\ \underbrace{\bigl(\forall k:\ \dom.w_{k} \le
\tau_{k}\bigr)}_{\text{\emph{domain} width within \emph{domain} tolerance}} ,
\end{equation} and its split halves the axis furthest past its own tolerance,
$\arg\max_{k} \dom.w_{k}/\tau_{k}$. The search starts from $[0,1]^{3}$, so
every midpoint is a dyadic rational, and it is exact in binary floating point
while the depth stays below the width of the significand. The implementation
does not assume that. It checks that the midpoint separates the endpoints, and
accepts the box when it does not.

The second, \textsc{Relaxed}, accepts on a \emph{codomain} width measured
against a \emph{domain} tolerance, $\forall d:
\overline{F}_{d}-\underline{F}_{d} \le \tau_{d}$. It adds the same error-box
condition and two degenerate cases. Comparing quantities that carry different
units makes the test loose, so it accepts earlier and reports a time of impact
further before the true one. \Cref{sec:results:relaxed} quantifies that trade,
which does not run one way.

Both tests also require $\dom.t_{\ell} > 0$. That refuses a contact reported at
exactly $t=0$, which would stall a solver using $\toi$ as a step length.
Measured against the exact roots the condition costs nothing: with it in place,
no collision is missed and no time of impact is late.

\subsection{The adaptive splitter, and what measuring it showed}
\label{sec:np:split}

\textsc{Relaxed} and the quad kernel split an axis several ways at once. The chosen axis is cut at $N=2$ interior points, giving three
children. Each point comes from a one-dimensional Gauss--Newton step against the
linearisation of $\Fdiff$ along that axis, clamped to within $0.45h$ of its
nominal position, where $h$ is the uniform spacing. The splitter is about
$1.6\times$ faster than a uniform grid, which needs five sub-intervals where
this uses three. Measuring \emph{why} produced a more useful result than the
speedup did.

First, the splitter barely adapts. The nominal point cancels exactly from the
Gauss--Newton step, so all $N$ splitters solve for the same least-squares point.
They differ only in which clamp window each one lands in, and $95.9\%$ of them
are pinned at their clamp. The clamp, and not the Newton step, is what makes
each child a bounded fraction of its parent. An unclamped version fails to
terminate.

Second, placement cannot affect correctness: the children tile the parent and
every child is still tested, so the splitter is purely a performance knob.

Third, and most informative, a principled alternative \emph{lost}. A
chord-based splitter aimed at the largest root-free region ran $2.7\times$
slower and isolated \emph{less} dead volume, $56.7\%$ against $61.1\%$. The
cause is a mismatch between what is dead and what is rejectable. On a typical
parent, $68.4\%$ of the volume is root-free \emph{pointwise}, taken as a union
over the three components of $\Fdiff$. Rejecting a child, however, requires a
\emph{single} component to exclude the origin over the whole child, and only
about $58\%$ of the axis is rejectable that way. A multi-way split does better
because different children may die by different components. A single
well-placed cut has no such advantage.

\subsection{Order independence}
\label{sec:np:order}

TightInclusion pops boxes from a priority queue ordered by ascending $t$ lower
bound, and returns the first box it accepts. Its answer is therefore the
accepted box of globally smallest $t$. \Cref{alg:np} keeps no such order. The
proposition below says that it needs none, and it is the result that licenses
everything in \cref{sec:parallel}.

\begin{proposition}[Order independence]
\label{prop:order}
Let $\mathcal{A}$ be the set of boxes that \op{Accept} accepts, over the tree
generated by a fixed \op{Split}. Suppose a search explores that tree in any
order, maintains $\toi = \min\{\dom.t_{\ell} : \dom \in \mathcal{A}
\text{ reached}\}$, and skips any box with $\dom.t_{\ell} \ge \toi$. Then on
termination $\toi = \min\{\dom.t_{\ell} : \dom \in \mathcal{A}\}$, independent of
the exploration order.
\end{proposition}

\begin{proof}
Let $m$ be the right-hand side and $\dom^{\ast}$ a box attaining it. Since $\toi$
only ever takes values $\dom.t_{\ell}$ for reached $\dom \in \mathcal{A}$, we
have $\toi \ge m$ throughout. For the converse, suppose $\dom^{\ast}$ is skipped.
It is skipped only because $\dom^{\ast}.t_{\ell} \ge \toi$ at that moment, and
$\toi \ge m = \dom^{\ast}.t_{\ell}$, so $\toi = m$ already. Otherwise
$\dom^{\ast}$ is reached and $\toi \le \dom^{\ast}.t_{\ell} = m$. The same
argument applies to every ancestor of $\dom^{\ast}$: an ancestor is pruned only
when $\toi$ has already reached a value at or below its $t$ lower bound, which is
at or below $m$. Hence $\toi = m$ in all cases.
\end{proof}

\begin{remark}[Agreement with the reference]
\label{rem:ti}
TightInclusion's answer is the same minimum, so the two agree exactly. Its queue
is ordered by ascending $t$ lower bound, so the first accepted box it pops has the
smallest $t$ lower bound among accepted boxes it \emph{reaches}; and a box it
does not reach was pruned by the same $t$ bound, which by the argument above is
already at or below that minimum. Hence its answer is
$\min\{\dom.t_{\ell} : \dom \in \mathcal{A}\}$, the same quantity
\cref{prop:order} gives.
\end{remark}

The hypotheses of \cref{prop:order} are not a formality. It fixes \op{Accept}
and \op{Split}, and so fixes $\mathcal{A}$. Two searches agree only when they
generate the same $\mathcal{A}$, and that requires the same acceptance
\emph{rule} together with the same computed values feeding it, down to the last
bit.

Our host kernel and TightInclusion meet that condition. The corner evaluation is
written in the reference's expression order precisely so that they do
(\cref{sec:prelim:err}). \Cref{sec:results:ref} confirms the consequence to
every printed digit on all twelve scene-phases, which is the sharpest check of
the proposition available to us.

Our \emph{device} kernel does not meet it. It evaluates the same formulas
through vector fused multiply-adds, and it rounds $\delta^{3}$ upward with a plain cube. A box decided in the last few units in the last place
can therefore fall differently, and its $\mathcal{A}$ is a slightly different
set. \Cref{prop:order} thus explains why the host reproduces the reference, and
says nothing about how closely the device will. \Cref{sec:results:ref} measures
that separately, and the answer is close without being identical.

What the proposition buys is freedom. The traversal may be depth-first,
breadth-first, or interleaved across queries; boxes from unrelated queries may
share a vector register or a thread block; work may migrate between processors
mid-search. None of that can change the answer. The two implementations in
\cref{sec:parallel} spend that freedom in different currencies.

\subsection{Quads}
\label{sec:np:quad}

Quad meshes are handled directly, without triangulation. They have their own
inclusion function (\cref{eq:fvq}), their own tolerance and error reductions,
their own host root finder and their own device kernel.

This path is \emph{not evaluated} in this paper. Every scene in
\cref{sec:results} is a triangle mesh, so no measurement reported here bears on
a quad, and the error constant of \cref{sec:prelim:err} for the bilinear form is
uncertified. What follows is a design, verified for correctness against a
densely sampled reference root and against the alternative factoring below. It
is not a measured result.

We considered extending the triangle kernel and rejected it. That kernel is
built end to end around a four-point query packed into one vector register, on a
barycentric domain whose validity test rejects $u+v \ge 1$. A vertex-quad query
has five points on an independent domain, so the load, the tolerances, the error
bound, the corner evaluation, the domain test and the split would all have had
to change at once.

Two details of the implementation matter. The first is the factoring. The
evaluation is written $\Fdiff = q(t) - \sum_{k} w_{k}(u,v)\,g_{k}(t)$, and not
as a blend in $(u,v)$ first, because the weights do not depend on $t$. A split
along $u$ or $v$ then holds the $t$-dependent frame fixed across every corner it
evaluates, so the frame is computed once for the whole split.

The second is the node order, which in \cref{eq:fvq} is lexicographic and not
cyclic. Winding the nodes around the quad swaps the last two. The solver then
searches a bilinear saddle stretched across the diagonal, and reports a
perfectly valid time of impact for a surface the caller did not mean. That
failure is silent, which is why we state the convention here and not in a
manual.

One further measurement from the quad path generalises. Hoisting the evaluation
of the split planes out of the loop over children is the textbook optimisation.
It is $1.3\times$ faster per isolated query and $1.35\times$ slower end to end.
The cause is that $\toi$ sharpens \emph{during} the loop. A child whose $t$
lower bound has since passed the running minimum is discarded unlooked-at, and
hoisting evaluates it anyway. When the bound is shared, a loop-invariant
computation is not invariant in the only sense that matters.
